\documentclass[10pt,a4paper]{amsart}
\usepackage[margin=19mm]{geometry}
\usepackage{amsmath,amssymb,mathtools}
\usepackage[scr=rsfs,cal=euler]{mathalfa}
\usepackage{xfrac}
\usepackage[unicode,hidelinks,bookmarks=false]{hyperref}
\newcommand{\ie}{i.e.\ }
\newcommand{\cf}{cf.\ }
\newcommand{\resp}{respectively\ }
\newcommand{\Subsection}{\subsection}
\providecommand{\nobreakdash}{\nobreak\hbox{-}}
\setlength{\emergencystretch}{2em}
\multlinegap0pt
\usepackage{amssymb}
\usepackage{mathtools}
\newtheorem{lemma}{Lemma}
\title{Complex Hyperbolic Immersions of Cheng--Yau Metrics on Thullen Domains}
\author{Mirel Caib{\u a}r, Andrea Loi}
\date{Source: arXiv:2608.21911v1 (CC BY 4.0)}
\begin{document}
\maketitle

\noindent\emph{Source and licence.} Complex Hyperbolic Immersions of Cheng--Yau Metrics on Thullen Domains. Version arXiv:2608.21911v1 (CC BY 4.0), distributed by arXiv under the Creative Commons Attribution 4.0 International licence, http://creativecommons.org/licenses/by/4.0/, as recorded in the arXiv licence metadata for this exact version and shown on the abstract page https://arxiv.org/abs/2608.21911v1. Full licence notice: \texttt{licenses/arxiv-2608.21911-CC-BY-4.0.txt}.\par\smallskip

\noindent\textit{Editorial scope.} Counted unit: the article's lemma lem:algebraic-uniformization (source0.tex lines 587--641). The article's complete original proof of it is delivered with it (source0.tex lines 643--706, one \texttt{proof} environment of 64 lines). No mathematics is written, repaired or re-derived: the statement and the proof are the article's own lines, in the article's own notation, and no retained line is substituted or re-flowed.  Retained uncounted as the source-local context the proof needs: lines 504--510; lines 514--516; lines 518--524; lines 564--571.  Nothing was omitted from any retained span, so the package carries no omission record.  No retained line contains a citation, so the wrapper carries no bibliography and no bibliography entry.  No retained line contains a figure environment, an image-inclusion command or drawing code, so no image is delivered and the package declares no asset dependency.  Editorial formatting only, in addition: an AMS \texttt{amsart} preamble that restates the article's own declarations and macros used by the retained lines, namely the environments \texttt{lemma} on separate counters, together with the standard packages those lines need.  The retained environments print in this document as Lemma 1 in order of appearance, whereas the article prints the same items with the numbering of its own full document.  The complete native source of the article as distributed in the arXiv e-print for this exact version is bundled unchanged as \texttt{sources/208285/source0.tex}.
\begin{equation}\label{eq:PQM}
 p=x\varphi'(x),
 \qquad
 \mathcal Q(p)=\mu p^2+(\mu+1)p+c_{\mathrm{ke}},
 \qquad
 \mathcal M(p)=c_{\mathrm{ke}}+\mu p.
\end{equation}

\begin{equation}\label{eq:first-integral}
 x e^{3\varphi}=p\,\mathcal Q(p).
\end{equation}

\begin{equation}\label{eq:p-separable}
 \frac{dp}{d\log x}
 =\frac{p\mathcal Q(p)}{\mathcal M(p)},
 \qquad
 \frac{d\log x}{dp}
 =\frac{\mathcal M(p)}{p\mathcal Q(p)}.
\end{equation}

\begin{equation}\label{eq:mu-c-r}
 \mu_r=\frac{3r(1-r)}{D_r},
 \qquad
 c_{\mathrm{ke},r}=\frac{\mu_r+2}{3}
 =\frac{(2-r)(1+r)}{3D_r},
 \qquad
 c_r=\frac{3r}{1+r}.
\end{equation}

\begin{lemma}[Algebraic parametrization]
\label{lem:algebraic-uniformization}
Let \(\varphi\) be the normalized solution for \(\mu=\mu_r\), and let
\[
 p=\xi\varphi'(\xi).
\]
Set
\begin{equation}\label{eq:b-r}
 b_r=\frac{3(1-r)}{2-r}.
\end{equation}
Then
\begin{equation}\label{eq:Q-factor}
 \mathcal Q_r(p)
 =c_{\mathrm{ke},r}(1+c_rp)(1+b_rp),
\end{equation}
and
\begin{equation}\label{eq:xi-p-factor}
 \xi=c_{\mathrm{ke},r}
 p(1+c_rp)^{-r}(1+b_rp)^{-(1-r)}.
\end{equation}

Define
\begin{equation}\label{eq:t-kappa}
 t=\frac{(b_r-c_r)p}{1+b_rp},
 \qquad
 \kappa_r=\frac{b_r}{b_r-c_r}
 =\frac{1-r^2}{1-2r},
\end{equation}
and
\begin{equation}\label{eq:Y-def}
 Y=\frac{b_r-c_r}{c_{\mathrm{ke},r}}\,\xi.
\end{equation}
Then
\begin{equation}\label{eq:Y-t}
 Y=t(1-t)^{-r}.
\end{equation}
If
\begin{equation}\label{eq:G-def}
 G_r(Y)=e^{-c_r\varphi(\xi(Y))},
\end{equation}
then
\begin{equation}\label{eq:G-t}
 G_r(Y)=(1-\kappa_rt)^{c_r}(1-t)^{-r},
\end{equation}
and
\begin{equation}\label{eq:derivative-cancellation}
 -\frac{dG_r}{dY}
 =A_r(1-\kappa_rt(Y))^{-\delta_r},
 \qquad
 A_r=\frac{r(2-r)}{1-2r},
 \qquad
 \delta_r=\frac{1-2r}{1+r}.
\end{equation}
In particular, \(A_r,\delta_r>0\) for \(0<r<1/2\).
\end{lemma}

\begin{proof}
Substituting \eqref{eq:mu-c-r} and \eqref{eq:b-r} into
\eqref{eq:PQM} gives \eqref{eq:Q-factor}. Moreover,
\[
 \frac{c_{\mathrm{ke},r}+\mu_rp}{p\mathcal Q_r(p)}
 =\frac1p-\frac{rc_r}{1+c_rp}
 -\frac{(1-r)b_r}{1+b_rp}.
\]
Since \(\varphi(0)=0\), evaluating the normalized equation at the origin
gives
\[
 c_{\mathrm{ke},r}\varphi'(0)=1.
\]
Thus
\[
 p=\xi\varphi'(\xi)\sim\frac{\xi}{c_{\mathrm{ke},r}},
 \qquad
 \frac{\xi}{p}\longrightarrow c_{\mathrm{ke},r}
 \quad\text{as }\xi\to0.
\]
Integrating \(d\log\xi/dp\) in \eqref{eq:p-separable} therefore yields
\eqref{eq:xi-p-factor}.

The change of variable \eqref{eq:t-kappa} satisfies
\begin{equation}\label{eq:linear-factors-t}
 1+b_rp=(1-\kappa_rt)^{-1},
 \qquad
 1+c_rp=\frac{1-t}{1-\kappa_rt}.
\end{equation}
Substitution into \eqref{eq:xi-p-factor} gives \eqref{eq:Y-t}; the powers
of \(1-\kappa_rt\) cancel because the exponents \(r\) and \(1-r\) add
to one.

Since \(Y(0)=0\) and
\[
 \left.\frac{dY}{dt}\right|_{t=0}=1,
\]
the relation \eqref{eq:Y-t} admits a real-analytic local inverse
\(t=t(Y)\) near the origin.

By \eqref{eq:first-integral},
\[
 e^{-c_r\varphi}
 =\left(\frac{\xi}{p\mathcal Q_r(p)}\right)^{c_r/3}.
\]
Since \(c_r/3=r/(1+r)\), equations \eqref{eq:Q-factor} and
\eqref{eq:linear-factors-t} give \eqref{eq:G-t}. Finally,
\[
 \frac{dY}{dt}
 =(1-t)^{-r-1}\bigl(1-(1-r)t\bigr),
\]
while
\[
 -\frac1{G_r}\frac{dG_r}{dt}
 =\frac{c_r\kappa_r}{1-\kappa_rt}-\frac{r}{1-t}
 =\frac{A_r(1-(1-r)t)}
 {(1-\kappa_rt)(1-t)}.
\]
Dividing by \(dY/dt\) and using
\[
 c_r-1=-\delta_r
\]
gives \eqref{eq:derivative-cancellation}.
\end{proof}

\end{document}
